\newcommand{\auditoriaid}{A03 -- Niño 3.4 frente a Niño 1+2 y métodos TOE}
% Preambulo comun de las auditorias del informe E1 (revision_e1/)
\documentclass[11pt,a4paper]{article}
\usepackage[utf8]{inputenc}
\usepackage[T1]{fontenc}
\usepackage[spanish,es-noshorthands]{babel}
\usepackage{lmodern}
\usepackage{microtype}
% Sin patrones de guionado en espanol en esta instalacion (ver A10): se
% desactiva el guionado automatico para evitar cortes con reglas del ingles.
\hyphenpenalty=10000
\exhyphenpenalty=10000
\usepackage[a4paper,margin=2.2cm]{geometry}
\usepackage{booktabs,array,longtable}
\usepackage{graphicx,caption,subcaption}
\usepackage{xcolor}
\usepackage{enumitem}
\usepackage{fancyhdr}
\usepackage{amsmath}
\usepackage{tikz}
\usetikzlibrary{arrows.meta,positioning}
\usepackage{natbib}
\usepackage{xurl}
\usepackage[hidelinks]{hyperref}
\urlstyle{tt}
\newcommand{\archivo}[1]{\url{#1}}
\newcolumntype{L}[1]{>{\raggedright\arraybackslash}p{#1}}
\definecolor{alta}{RGB}{180,30,30}
\definecolor{media}{RGB}{200,120,0}
\definecolor{baja}{RGB}{60,110,160}
\definecolor{cajafondo}{RGB}{245,247,250}
\newcommand{\sevA}{\textcolor{alta}{\textbf{Alta}}}
\newcommand{\sevM}{\textcolor{media}{\textbf{Media}}}
\newcommand{\sevB}{\textcolor{baja}{\textbf{Baja}}}
\setlength{\parskip}{0.35em}
\setlength{\emergencystretch}{3em}
\setlength{\parindent}{0pt}
\pagestyle{fancy}
\fancyhf{}
\fancyhead[L]{\small Auditoría del Informe del Entregable~1 -- \auditoriaid}
\fancyhead[R]{\small \thepage}
\renewcommand{\headrulewidth}{0.4pt}
% Entorno de hallazgos: ID | Severidad | Ubicacion | Hallazgo y evidencia | Correccion
\newenvironment{hallazgos}{%
\small
\begin{longtable}{@{}L{1.15cm}L{1.25cm}L{1.9cm}L{6.0cm}L{\dimexpr\textwidth-10.3cm-10.6\tabcolsep\relax}@{}}
\toprule
\textbf{ID} & \textbf{Sev.} & \textbf{Ubicación} & \textbf{Hallazgo y evidencia} & \textbf{Corrección aplicada} \\
\midrule\endfirsthead
\toprule
\textbf{ID} & \textbf{Sev.} & \textbf{Ubicación} & \textbf{Hallazgo y evidencia} & \textbf{Corrección aplicada} \\
\midrule\endhead
\bottomrule\endfoot
}{\end{longtable}}
% Macros compartidas con el informe revisado (las secciones propuestas las usan)
% Macros compartidas entre el informe revisado y las auditorias.
% Cifras verificadas contra los archivos del repositorio (ver A01/A04/A07):
\newcommand{\nUniverso}{102}   % source_id CMIP6 con tos/Omon en ESGF (00b, model_availability_priority.csv)
\newcommand{\nSeed}{41}        % publican los cuatro experimentos (config/models_seed_cmip6.csv)
\newcommand{\nSel}{40}         % completos + QC PASS (models_inventory_final.csv)
\newcommand{\nComb}{160}       % combinaciones modelo-experimento PASS (qc_report.csv)
\newcommand{\nNoEnc}{44}       % config/models_no_encontrado_44.csv
\newcommand{\nParcial}{18}     % 102 - 40 - 44
% Fecha de la portada: fija (no \today), editable en un solo lugar.
\newcommand{\fechaentrega}{28 de septiembre de 2026}
% Estado de codigo que documenta el E1 (antes de los cambios del E2)
\newcommand{\commitEone}{\texttt{8ee6184}}

% En las auditorias el texto propuesto se muestra fuera del informe: las
% referencias cruzadas se imprimen con el nombre de la seccion destino.
\makeatletter
\def\etq#1#2{\expandafter\def\csname etq@#1\endcsname{#2}}
\etq{sec:resumen}{Resumen}
\etq{sec:alcance}{Alcance}
\etq{sec:marco}{Revisión científica}
\etq{sec:antecedentes}{Antecedentes}
\etq{sec:referencias}{Literatura sobre TOE}
\etq{tab:referencias}{bibliografía TOE}
\etq{sec:similitudes}{Niño~3.4 frente a Niño~1+2}
\etq{sec:toe_metodos}{Elección de los métodos de TOE}
\etq{tab:metodos_toe}{métodos TOE}
\etq{eq:szabo_signal}{ecuación señal M1}
\etq{eq:szabo_V}{ecuación V(t)}
\etq{eq:szabo_T}{ecuación T(t)}
\etq{eq:szabo_toe}{ecuación TOE M1}
\etq{eq:tod}{ecuación ToD}
\etq{sec:datos}{Datos}
\etq{fig:cmip6_contribution}{contribución institucional}
\etq{sec:tabla_modelos}{Modelos seleccionados}
\etq{tab:modelos}{modelos}
\etq{sec:desviacion}{Selección de modelos}
\etq{sec:dominios_periodos}{Dominios y periodos}
\etq{fig:region_nino}{regiones Niño}
\etq{sec:arquitectura}{Arquitectura del pipeline}
\etq{fig:flujograma}{Flujograma}
\etq{sec:scripts}{Descripción de los pasos}
\etq{sec:paso00b}{Paso 00b}
\etq{sec:paso01}{Paso 01}
\etq{sec:paso02}{Paso 02}
\etq{sec:paso04}{Paso 04}
\etq{sec:paso05}{Paso 05}
\etq{sec:paso06}{Paso 06}
\etq{sec:paso07}{Paso 07}
\etq{sec:reporte}{Reporte de estado}
\etq{sec:resultados}{Resultados}
\etq{fig:qc}{matriz de estado}
\etq{tab:cuentas}{conteo de modelos}
\etq{fig:mapas2d}{mapas 2D}
\etq{fig:series}{series}
\etq{fig:boxplot}{boxplots}
\etq{sec:roadmap}{Próximos pasos}
\etq{tab:roadmap}{próximos pasos}
\etq{sec:conclusiones}{Conclusiones}
\etq{sec:recomendaciones}{Recomendaciones}
\makeatother

\makeatletter
\AtBeginDocument{\renewcommand{\ref}[1]{\@ifundefined{etq@#1}{\textit{[#1]}}{\textit{\csname etq@#1\endcsname}}}}
\makeatother
% Texto propuesto: secciones degradadas un nivel y en caja
\newenvironment{propuesto}{%
  \par\medskip\noindent\textcolor{baja}{\rule{\textwidth}{0.8pt}}\par
  \noindent{\small\textbf{Inicio del texto propuesto para el informe revisado}}\par\medskip
  \begingroup
  \let\section\subsection\let\subsection\subsubsection\let\subsubsection\paragraph
}{%
  \endgroup
  \par\noindent{\small\textbf{Fin del texto propuesto}}\par
  \noindent\textcolor{baja}{\rule{\textwidth}{0.8pt}}\par\medskip}
% Recuadro de actualizacion posterior (decisiones del autor)
\newcommand{\actualizacion}[1]{\par\medskip\noindent\fcolorbox{media}{media!8}{\parbox{\dimexpr\linewidth-2\fboxsep-2\fboxrule\relax}{\small\textbf{Actualización del 29-09-2026 (decisiones del autor).} #1}}\par\medskip}

\begin{document}

\begin{center}
{\Large\bfseries Auditoría A03: Niño~3.4 frente a Niño~1+2 y métodos de TOE}\\[0.4em]
{\normalsize Informe del Entregable~1, Secciones 3.3 y 3.4 (\texttt{informe\_e1\_smnv3.tex}, commit \texttt{330fd2e})}\\[0.2em]
{\small Revisión: \fechaentrega}
\end{center}

\section{Objeto}

Se auditan la Sección~3.3 (``Similitudes y diferencias entre Niño~3.4 y Niño~1+2'') y la Sección~3.4 (elección de los métodos de TOE, con las ecuaciones de M1 y ToD). Es la parte del informe con más afirmaciones cuantitativas atribuidas a terceros y la que fija el método del Entregable~4, por lo que cada número y cada ecuación se contrastó con el texto completo de su fuente. Además, se verificó con datos propios (ERSSTv5) la premisa física de la hipótesis de asimetría Niño~1+2/Niño~3.4.

\section{Verificación de ecuaciones de M1 contra \citet{szabo2016}}

\begin{center}\small
\begin{tabular}{@{}L{2.2cm}L{5.4cm}L{\dimexpr\textwidth-7.6cm-4\tabcolsep\relax}@{}}
\toprule
Elemento & Capítulo de Szabó (Ecs.~12.1--12.9) & Informe original \\
\midrule
Señal & Polinomio de 4.º orden; $X=\Delta T+\bar T+\varepsilon$ (12.1) & Correcto. \\
$\bar T$ & Promedio del ajuste en 1971--2000 & Da 1981--2014 sin aclarar que es la adaptación del proyecto. \\
$V(t)$ & $\mathrm{var}_{m,s,t}$ de $\varepsilon$ en los 30 años previos, recalculada cada año; media móvil de 10 años & Omite el suavizado; la ventana $[t-30,t]$ tiene 31 años. \\
$M(t)$, $S(t)$ & (12.4) media sobre escenarios de $\mathrm{var}_m$; (12.5) $\mathrm{var}_s$ de la media multimodelo & Solo en prosa y de forma imprecisa. \\
$T(t)$ & $V+M+S$ (12.6), igual que HS2009 Ec.~7 & Presentado como rasgo propio de Szabó frente a la ``raíz de la suma de cuadrados''. \\
$G(t)$ & Media sobre $m,s$ de $\Delta T$ (12.7) & Usada en STN y TOE sin definirse. \\
STN & $G/\sqrt{T}$ (12.8), sin test de significancia & Correcto. \\
TOE & $G/\sqrt{V}\geq\kappa$, ``usualmente $+1$ o $-1$'' (12.9) & $\kappa\in\{1,2\}$ ``siguiendo la misma referencia'': el 2 no es de Szabó. \\
\bottomrule
\end{tabular}
\end{center}

La ecuación del ToD coincide con \citet{gopika2024} (Ecs.~3--4; grados de libertad efectivos según Bayley y Hammersley, 1946; prueba $t$ unilateral), pero el informe no define $W$ (índice de calentamiento tropical, TSM media 30°S--30°N) ni la regla ``último año desde el que la significancia se mantiene''.

\section{Verificación con datos propios}

La hipótesis final de la Sección~3.3 se apoya, vía \citet{chadwick2019}, en que mayor variabilidad implica un TOE más tardío. Se calculó la variabilidad observada de ambas cajas con ERSSTv5 (\archivo{data/processed/masked/ersstv5_region.nc}, 1981--2014; media ponderada por $\cos\phi$, anomalías respecto del ciclo anual, sin tendencia lineal):

\begin{center}\small
\begin{tabular}{@{}lrrrr@{}}
\toprule
Caja & Puntos de grilla & TSM media (°C) & DE anomalía mensual (°C) & DE media anual (°C) \\
\midrule
Niño~3.4 & 130 & 26{,}99 & 0{,}86 & 0{,}61 \\
Niño~1+2 & 35 & 23{,}32 & 1{,}08 & 0{,}86 \\
\bottomrule
\end{tabular}
\end{center}

Niño~1+2 tiene un ruido interanual entre 25 y 40\,\% mayor. La premisa se confirma y se incorpora al texto como evidencia propia.

\section{Hallazgos}

\begin{hallazgos}
A03-1 & \sevA & 3.3, Gopika & ``coincide geográficamente con Niño~1+2/3.4''. La caja EEP de esta literatura (Y22 y \citealp{zhong2023}: 2{,}5°S--2{,}5°N, 180°--90°O) se superpone a Niño~3.4 y no incluye Niño~1+2 (90°--80°O, 10°S--0°). & Se describe la superposición real. \\
A03-2 & \sevM & 3.3, Gopika & ``91\,\% de los modelos''. El artículo dice ``$\sim$90\,\%'' en el resumen y ``over 90\,\%'' en el texto. Omite además el resultado más cercano a Niño~1+2: calentamiento atenuado del Pacífico sudeste, robusto en observaciones y en $\sim$80\,\% de los modelos. & ``más del 90\,\%''; se agrega el resultado del Pacífico sudeste. \\
A03-3 & \sevM & 3.3, Gopika & ``esto sustenta por qué el TdR exige \ldots\ determinar los sesgos de la TSM en el Pacífico tropical central y oriental''. El TdR, actividad (b), pide sesgos sin especificar región. & ``hay que evaluar el sesgo \ldots\ tarea del Entregable~2''. \\
A03-4 & \sevM & 3.3, Zhong & ``77--80\,\% en CESM1-LE'' mezcla dos cifras: $\sim$80\,\% para la TSM del ENOS y $\sim$77\,\% para la precipitación. No da el 21\,\% del estado medio en el EEP, que es el contraste relevante. & Cifras separadas; se agrega el 21\,\%. \\
A03-5 & \sevM & 3.3, Zhong & ``exactamente la separación que el TdR exige entre \ldots\ estado medio (E3) y partición de incertidumbre (E4)''. La actividad (c) del TdR incluye estado medio \emph{y} variabilidad; la correspondencia es forzada. & Se dice que deben analizarse por separado en E3 y E4. \\
A03-6 & \sevA & 3.3, Larson y Poul & ``documentan \ldots\ el mismo mecanismo físico: la circulación forzada por el viento \ldots\ amplifica el ruido''. Poul et al. atribuyen el retraso a una surgencia intensificada que \emph{reduce la señal} (y reduce o no cambia la variabilidad), no a más ruido. & ``por dos vías distintas'': más ruido (Larson), menos señal (Poul). \\
A03-7 & \sevA & 3.3, Larson & ``en ocasiones impidiendo que la señal emerja del todo dentro del siglo XXI''. El análisis de Larson y McMonigal usa ensambles \emph{históricos} de CESM2 (1900--2014); la falta de emergencia es dentro de ese periodo. & ``para el periodo histórico (1900--2014)''. \\
A03-8 & \sevM & 3.3, Larson & ``la firma física de la Corriente de Humboldt frente a Niño~1+2'' se presenta como dicho por el artículo; el artículo habla de surgencia en la frontera oriental del Pacífico, sin nombrar Humboldt. & La mención de Humboldt se presenta como lectura del equipo. \\
A03-9 & \sevB & 3.3, Keller & ``atribuido a la alta variabilidad natural asociada a ENOS'': la frase del artículo sobre variabilidad del ENOS se refiere a Friedrich et al. (2012) y al Pacífico ecuatorial oriental, no a la zona de Perú y Chile. & Se omite la atribución y se da el dato verificable (TOE de TSM de 45--90 años; $>$30 años para pCO$_2$ y pH frente a Perú y Chile). \\
A03-10 & \sevM & 3.3, Lois & ``no colapse artificialmente la dispersión de Índice~C e Índice~E''. El commit \texttt{38dd573} anunció que se quitaban las menciones de los Índices C y E, pero esta quedó; el Entregable~2 usa ONI/RONI e ICEN. & ``dispersión de la respuesta de la TSM en ambas cajas''. \\
A03-11 & \sevM & 3.3, cierre & ``cuatro líneas de evidencia independientes \ldots\ convergen''. Zhong trata regímenes de incertidumbre, no el orden de emergencia; Chadwick trata precipitación en Chile. Además, el párrafo repite las cuatro citas (comentario \texttt{\% REVISAR}). & Cierre de una frase como ``hipótesis de trabajo''; la analogía de Chadwick se respalda con ERSSTv5. \\
A03-12 & \sevB & 3.3, título & ``Similitudes y diferencias entre Niño~3.4 y Niño~1+2'' no describe el contenido, que es una lectura de la bibliografía. & ``Implicancias de la bibliografía para Niño~3.4 y Niño~1+2''. \\
A03-13 & \sevA & 3.4, Ec. TOE & ``$\kappa\in\{1,2\}$ \ldots\ siguiendo la misma referencia de \citet{szabo2016}''. Szabó usa $\pm1$; el umbral 2 proviene de \citet{hawkins2012}. & Se atribuye cada umbral a su fuente. \\
A03-14 & \sevA & 3.4, Ec. $T(t)$ & ``suma aditiva \ldots\ (no la raíz de la suma de cuadrados usada por otros métodos)''. Es la misma magnitud (ver A02-3), heredada de HS2009. & Se elimina el contraste. \\
A03-15 & \sevM & 3.4, M1 & $G(t)$ no se define; $M$ y $S$ no tienen ecuación; falta el suavizado de 10 años de $V$; la ventana $[t-30,t]$ tiene 31 años; ``mismo periodo de referencia'' que Bruno (1981--2010 frente a 1981--2014). & Se definen $G$, $M$ y $S$; ventana $[t-29,t]$; suavizado; ``periodo similar''. \\
A03-16 & \sevM & 3.4, ToD & No se definen $W$ ni la regla del ``último año''. $W$ es la TSM media 30°S--30°N, un dominio mayor que la ventana de descarga del E1 (20°S--20°N); la ampliación ya se hizo en el E2 (\texttt{79a9786}, por RONI). & Se definen $W$ y la regla. \textit{Actualización 29-09: la ventana de descarga vigente (0°--360°, 30°S--30°N) ya cubre $W$, y así se indica.} \\
A03-17 & \sevM & 3.4, criterios & ``(ii) ToD ya fue validado \ldots\ exactamente sobre la caja EEP que coincide con Niño~3.4/1+2'' (ver A03-1); los métodos descartados no se citan. & Redacción corregida y citas \citet{john2023}, \citet{barrow2019}, \citet{ryu2024}. \\
A03-18 & \sevB & 3.4 & Comentario \texttt{\% REVISAR} sobre si conservar las ecuaciones en el E1. Nombres inconsistentes: ``M2 complementario ToD'' frente a ``ToD (Gopika)''. & Se conservan y se explica por qué (fijar el método antes de calcularlo); nombre único ``M2 (ToD)''. \\
\end{hallazgos}

\section{Extensión}

Original: Sección~3.3, 657 palabras; Sección~3.4, 687 (1344 en total). Propuesto: 1126 
palabras. La reducción viene de eliminar el párrafo de cierre que repetía las cuatro citas y las afirmaciones no sustentadas, aun después de agregar las definiciones que faltaban ($G$, $M$, $S$, $W$) y la verificación con ERSSTv5.

\section{Texto propuesto}

\begin{propuesto}
\subsection{Implicancias de la bibliografía para Niño~3.4 y Niño~1+2}
\label{sec:similitudes}

Siete referencias del Cuadro~\ref{tab:referencias} merecen una lectura específica frente al enfoque de este proyecto.

\citet{gopika2024} no solo aportan el método ToD, adoptado como segunda estimación del TOE: analizan 69 simulaciones CMIP5/6 y cuatro productos observacionales en todo el trópico (30°S--30°N) y encuentran que el Pacífico ecuatorial oriental (EEP, por sus siglas en inglés) es la excepción principal. Allí el calentamiento es detectable en más del 90\,\% de los modelos, pero solo en uno de los cuatro productos observacionales, porque los modelos simulan un calentamiento más intenso que el observado. La caja EEP de esta literatura (2{,}5°S--2{,}5°N, 180°--90°O en \citealp{zhong2023}) se superpone a Niño~3.4 y no incluye Niño~1+2. Los mismos autores detectan, en observaciones y en cerca del 80\,\% de los modelos, un calentamiento del Pacífico sudeste más débil que el promedio tropical. Ambos resultados indican que, antes de proyectar, hay que evaluar el sesgo de los modelos en el Pacífico tropical oriental, tarea del Entregable~2.

\citet{zhong2023} establecen una distinción que este proyecto adopta como principio de diseño. Con tres grandes ensambles y CMIP6, muestran que en el EEP la variabilidad interna explica menos de la mitad de la incertidumbre del TOE de la TSM media anual (21\,\% en CESM1-LE), de modo que dominan las diferencias entre modelos. En cambio, explica cerca del 80\,\% de la incertidumbre del TOE de la TSM asociada al ENOS (77\,\% para la precipitación) en CESM1-LE, el único ensamble con suficientes miembros con señal emergente. El estado medio y la variabilidad del ENOS no comparten, por tanto, el mismo régimen de incertidumbre, y deben analizarse por separado en los Entregables~3 y~4.

\citet{larson2025} y \citet{poul2025} documentan, por dos vías distintas, que la dinámica oceánica costera retrasa el TOE de la TSM. En dos ensambles de CESM2 para el periodo histórico (1900--2014), \citet{larson2025} muestran que la variabilidad interna de la circulación forzada por el viento aumenta el ruido y retrasa el TOE en casi todo el océano, más de una década en promedio sobre las latitudes bajas. Además, a lo largo de la frontera oriental del Pacífico, la ausencia de señal emergente coincide con regiones de surgencia forzada por viento. En el Mar del Norte y el Báltico, \citet{poul2025} atribuyen los TOE más tardíos a una surgencia más intensa, que reduce la señal de calentamiento. Frente a Niño~1+2, dominado por la surgencia costera del sistema de la Corriente de Humboldt, ambos resultados apuntan a una emergencia más tardía o más incierta.

\citet{keller2014} reportan, con 17 modelos del sistema terrestre, un TOE de la TSM de 45 a 90 años, mucho mayor que el de las variables biogeoquímicas, e identifican la zona de surgencia de Perú y Chile como una de las regiones con TOE localmente elevado ($>$30 años) para pCO$_2$ y pH. Esta tercera línea, de naturaleza biogeoquímica, es independiente de las anteriores y apunta en el mismo sentido.

\citet{chadwick2019} desarrollan un TOE local para tres cuencas chilenas (Limarí, Maipo y Maule) y concluyen que el TOE de la precipitación es más tardío cuanto mayor es el coeficiente de variación local, con diferencias de más de 40 años según el percentil de tendencia de los modelos. Por analogía, la mayor variabilidad interanual observada en Niño~1+2 anticipa una emergencia más tardía que en Niño~3.4: con ERSSTv5 (1981--2014), la desviación estándar de las anomalías mensuales es de 1{,}08\,°C en Niño~1+2 frente a 0{,}86\,°C en Niño~3.4, y la de las medias anuales, de 0{,}86\,°C frente a 0{,}61\,°C.

\citet{lois2026} aportan un criterio de selección de modelos distinto de la validación física y estadística de \citet{bruno2023}: en lugar de filtrar por desempeño promedio, eligen un subconjunto de modelos CMIP6 que abarca respuestas de circulación y sensibilidades climáticas cualitativamente distintas. Aplicado a este proyecto, el principio sugiere verificar en los próximos entregables que la selección final de modelos no reduzca artificialmente la dispersión de la respuesta de la TSM en ambas cajas, para no subestimar la incertidumbre de modelo ni el rango del TOE en Niño~1+2.

En conjunto, estas líneas de evidencia sustentan una hipótesis de trabajo contrastable en el Entregable~4: \textbf{Niño~1+2 emergerá más tarde, o con mayor incertidumbre, que Niño~3.4}.

\subsection{Elección de los dos métodos de \emph{Time of Emergence}}
\label{sec:toe_metodos}

El TdR prevé para el Entregable~4 el cálculo de la razón señal-ruido y del TOE ``empleando metodologías revisadas''. El equipo técnico adopta como método M1 la formulación de \citet{szabo2016} y la complementa con un segundo método independiente, el \emph{Time of Detection} (ToD) de \citet{gopika2024}. Las ecuaciones se fijan desde este entregable para que el método quede definido antes de calcularlo; el cálculo corresponde al Entregable~4.

\subsubsection{Método M1 (Szabó y Szépszó)}

La señal se estima ajustando un polinomio de cuarto orden a la serie anual o estacional de TSM de cada modelo $m$ y escenario $s$:
\begin{equation}
\label{eq:szabo_signal}
X_{m,s}(t) = \Delta T_{m,s}(t) + \bar{T} + \varepsilon_{m,s}(t),
\end{equation}
donde $\bar{T}$ es el promedio del ajuste en el periodo de referencia (1971--2000 en \citealp{szabo2016}; 1981--2014 en este proyecto, Sección~\ref{sec:dominios_periodos}), $\Delta T_{m,s}(t)$ es la señal de cambio y $\varepsilon_{m,s}(t)$ es el residuo. A diferencia de \citet{hawkins2009}, la variabilidad interna no es constante: es la varianza de los residuos de todos los modelos y escenarios en los 30 años previos, recalculada cada año y suavizada con una media móvil de 10 años:
\begin{equation}
\label{eq:szabo_V}
V(t) = \mathrm{var}_{m,s,\tau}\big(\varepsilon_{m,s}(\tau)\ \big|\ \tau \in [t-29,\,t]\big).
\end{equation}
La incertidumbre de modelo es la varianza entre modelos promediada sobre los escenarios, $M(t)=\frac{1}{N_s}\sum_s \mathrm{var}_m\,\Delta T_{m,s}(t)$; la de escenario es la varianza entre escenarios de la media multimodelo, $S(t)=\mathrm{var}_s\big(\frac{1}{N_m}\sum_m \Delta T_{m,s}(t)\big)$; y, como en \citet{hawkins2009}, la incertidumbre total es su suma:
\begin{equation}
\label{eq:szabo_T}
T(t) = V(t) + M(t) + S(t).
\end{equation}
Con la señal media $G(t)=\frac{1}{N_m N_s}\sum_{m,s}\Delta T_{m,s}(t)$, la razón señal-ruido usa la incertidumbre total, $\mathrm{STN}(t) = G(t)/\sqrt{T(t)}$, mientras que el TOE usa solo la variabilidad interna como ruido:
\begin{equation}
\label{eq:szabo_toe}
\mathrm{TOE} = \min\,\Big\{t : \frac{G(t)}{\sqrt{V(t)}} \geq \kappa\Big\}.
\end{equation}
Se adoptan dos umbrales: $\kappa=1$, el de \citet{szabo2016}, y $\kappa=2$, el umbral más exigente de \citet{hawkins2012}. Separar el ruido de la STN (total) del ruido del TOE (solo interno) es el rasgo distintivo del método. \citet{bruno2023} ya lo aplicó con CMIP6 sobre Sudamérica, con un periodo de referencia similar, lo que reduce el riesgo de implementación.

\subsubsection{Método M2 (ToD, Gopika)}

Como segundo método se adopta el \emph{Time of Detection} de \citet{gopika2024}. En lugar de un umbral fijo de señal-ruido, se evalúa la significancia del coeficiente de una regresión sobre una ventana que crece desde un año inicial $t_0$ hasta $T$:
\begin{equation}
\label{eq:tod}
Y_T(t) = a_T\,W_T(t) + b_T + \varepsilon_T(t), \qquad n_{\text{eff}} = \frac{T-t_0+1}{\sum_{\tau=1}^{n/2} r^2(\tau)},
\end{equation}
donde $Y$ es la TSM de la caja, $W$ es el índice de calentamiento tropical (TSM media 30°S--30°N), $r(\tau)$ es la autocorrelación con desfase $\tau$ y $n_{\text{eff}}$ son los grados de libertad efectivos de la prueba $t$ unilateral sobre $a_T$. El ToD es el último año a partir del cual el coeficiente se vuelve significativo al 95\,\% y se mantiene significativo hasta el final de la serie. No depende de un umbral $\kappa$ sino del nivel de significancia; según los autores, el 95\,\% equivale aproximadamente a una razón señal-ruido de 1{,}6. El índice $W$ requiere la TSM de toda la franja tropical, que queda cubierta por la ventana de descarga del \emph{pipeline} (0°--360°, 30°S--30°N; Sección~\ref{sec:dominios_periodos}).

\begin{table}[htbp]
\caption{Métodos de TOE adoptados.}
\label{tab:metodos_toe}
\small
\begin{tabular}{@{}L{2.3cm}L{5.6cm}L{\dimexpr\textwidth-2.3cm-5.6cm-4\tabcolsep\relax}@{}}
\toprule
\textbf{Método} & \textbf{Ruido de referencia} & \textbf{Criterio de TOE} \\
\midrule
M1 (Szabó) & Variabilidad interna $V(t)$, ventana móvil de 30 años & $G(t)/\sqrt{V(t)} \geq \kappa$, con $\kappa=1$ o $2$ \\
M2 (ToD, Gopika) & Residuo de la regresión, con grados de libertad efectivos & Significancia del coeficiente al 95\,\%, sostenida hasta el final \\
\bottomrule
\end{tabular}
\end{table}

La elección de estos dos métodos responde a tres criterios: (i) M1 cumple directamente lo previsto para el Entregable~4, pues calcula la partición de incertidumbre, la razón señal-ruido y el TOE; (ii) el ToD fue aplicado por \citet{gopika2024} a la TSM tropical, incluido el Pacífico ecuatorial oriental, a diferencia de otros métodos revisados que se aplicaron a ríos australianos \citep{john2023}, a las praderas canadienses \citep{barrow2019} o a Corea del Sur \citep{ryu2024}; y (iii) ambos métodos son conceptualmente independientes (uno se basa en la razón señal-ruido con partición de incertidumbre; el otro, en la significancia estadística de una tendencia), lo que permite reportar un rango de fechas de emergencia en vez de un único número.

\end{propuesto}

\bibliographystyle{plainnat}
\bibliography{../informe/references_rev}
\end{document}
